\documentclass[11pt,a4paper]{article}

% Guida personale integrata: non è un capitolo definitivo della tesi.
\usepackage[utf8]{inputenc}
\usepackage[T1]{fontenc}
\usepackage[italian]{babel}
\usepackage{lmodern}
\usepackage{amsmath,amssymb,mathtools}
\usepackage[a4paper,margin=2.45cm]{geometry}
\usepackage{microtype}
\usepackage{enumitem}
\usepackage{booktabs,tabularx,array}
\usepackage{xcolor}
\usepackage{tikz}
\usetikzlibrary{arrows.meta,positioning,calc,shapes.geometric}
\usepackage{hyperref}

\definecolor{mainblue}{RGB}{48,94,151}
\definecolor{mainred}{RGB}{180,70,65}
\definecolor{maingreen}{RGB}{65,130,105}
\definecolor{mainorange}{RGB}{218,145,55}
\definecolor{softgray}{RGB}{245,246,248}

\hypersetup{colorlinks=true,linkcolor=mainblue,urlcolor=mainblue}
\setlength{\parindent}{0pt}
\setlength{\parskip}{0.55em}
\setlist{itemsep=0.2em,topsep=0.4em}

\tikzset{
  axis/.style={->,thick,gray!75},
  guide/.style={dashed,gray!65},
  data/.style={fill=mainblue!75,draw=none},
  extreme/.style={fill=mainred!80,draw=none},
  flowbox/.style={rounded corners,draw=mainblue!70,fill=mainblue!6,
                  minimum height=1.15cm,align=center,inner sep=6pt},
  flowarrow/.style={-{Latex[length=2.5mm]},thick,mainblue!75}
}

\newcommand{\Prob}{\mathbb{P}}
\newcommand{\E}{\mathbb{E}}
\newcommand{\one}{\boldsymbol{1}}
\newcommand{\ind}{\mathbf{1}}
\newcommand{\Idea}{\paragraph{Idea.}}
\newcommand{\Esempio}{\paragraph{Esempio.}}
\newcommand{\Ricorda}{\paragraph{Da ricordare.}}
\newcommand{\Attenzione}{\paragraph{Attenzione.}}
\newcommand{\Salta}{\paragraph{Cosa puoi saltare alla prima lettura.}}

\title{Tail dependence: una visione d'insieme\\
\large Guida personale integrata ai paper di Schmidt, Hua--Joe e Durante et al.}
\author{Appunti personali di Massimo}
\date{}

\begin{document}
\maketitle
\tableofcontents
\newpage

\section{La domanda generale: che cosa succede quando i rischi diventano estremi?}

Questa guida non riassume i tre paper uno dopo l'altro. Prova invece a
costruire un'unica storia. Il punto di partenza è finanziario: due asset possono
avere una relazione abbastanza tranquilla nei giorni normali e comportarsi in
modo molto diverso durante una crisi. Se le perdite grandi arrivano insieme,
la diversificazione può funzionare molto meno del previsto.

La correlazione lineare descrive soprattutto il comportamento medio. È utile,
ma non basta a rispondere a domande come:
\begin{itemize}
  \item se un titolo entra nel suo 1\% peggiore, quanto è probabile che anche
  l'altro sia nel proprio 1\% peggiore?
  \item due modelli con dipendenza asintotica nulla sono davvero equivalenti
  nelle code osservabili?
  \item come cambia la risposta passando dal limite teorico a una soglia
  finita, per esempio 5\%?
\end{itemize}

I tre paper costruiscono tre livelli della risposta:
\begin{enumerate}
  \item Schmidt fornisce il linguaggio operativo del tail-dependence
  coefficient (TDC), cioè il comportamento limite;
  \item Hua--Joe raffinano questo linguaggio con il tail order $\kappa$, cioè
  la velocità con cui una probabilità congiunta estrema scende a zero;
  \item Durante et al. mettono al centro la tail concentration function (TCF),
  cioè l'intera curva della concentrazione di coda a soglie finite.
\end{enumerate}

\begin{figure}[ht]
\centering
\begin{tikzpicture}[x=0.41cm,y=0.41cm,font=\small]
  \begin{scope}
    \draw[axis] (-4.2,0)--(4.4,0) node[right] {$X_1$};
    \draw[axis] (0,-4.2)--(0,4.4) node[above] {$X_2$};
    \draw[guide] (-1.45,-4)--(-1.45,4);
    \draw[guide] (-4,-1.45)--(4,-1.45);
    \node at (0,5.15) {relazione centrale simile};
    \foreach \x/\y in {-3.4/-0.2,-2.7/-1.0,-2.1/-1.3,-1.7/-2.6,-1.0/-0.8,-0.4/-1.7,0.3/0.2,0.9/1.0,1.4/1.7,2.2/1.1,2.8/2.6,3.3/2.0,-3.0/1.8,-2.0/2.5,1.8/-1.0,2.9/-0.8}{\fill[data] (\x,\y) circle (1.6pt);}
    \node[align=center] at (0,-5.0) {pochi estremi\\congiunti};
  \end{scope}
  \begin{scope}[xshift=8cm]
    \draw[axis] (-4.2,0)--(4.4,0) node[right] {$X_1$};
    \draw[axis] (0,-4.2)--(0,4.4) node[above] {$X_2$};
    \draw[guide] (-1.45,-4)--(-1.45,4);
    \draw[guide] (-4,-1.45)--(4,-1.45);
    \fill[mainred!10] (-4,-4) rectangle (-1.45,-1.45);
    \node at (0,5.15) {relazione centrale simile};
    \foreach \x/\y in {-3.4/-0.2,-2.7/-1.0,-2.1/-1.3,-1.0/-0.8,-0.4/-1.7,0.3/0.2,0.9/1.0,1.4/1.7,2.2/1.1,2.8/2.6,3.3/2.0,-3.0/1.8,-2.0/2.5,1.8/-1.0,2.9/-0.8}{\fill[data] (\x,\y) circle (1.6pt);}
    \foreach \x/\y in {-3.4/-3.1,-2.9/-2.6,-2.5/-3.2,-2.0/-2.8,-1.6/-3.5,-1.2/-2.9}{\fill[extreme] (\x,\y) circle (1.8pt);}
    \node[align=center] at (0,-5.0) {cluster di estremi\\nella lower tail};
  \end{scope}
\end{tikzpicture}
\caption{Due nuvole possono sembrare simili al centro e differire proprio nella regione che conta durante una crisi. Diagramma concettuale, non dati reali.}
\label{fig:corr-tail}
\end{figure}

\Ricorda Il problema della tail dependence non è sostituire la correlazione in
ogni situazione. È aggiungere una lente specifica per le regioni estreme.

\section{Dai rendimenti ai percentili: perché entra in gioco la copula}

\subsection{Osservazione, marginale e dipendenza}

Consideriamo due rendimenti $X_1$ e $X_2$. Una singola osservazione bivariata è
una coppia $(x_1,x_2)$ rilevata nello stesso giorno. La distribuzione marginale
$F_1$ descrive $X_1$ preso da solo; $F_2$ descrive $X_2$ preso da solo.

Trasformiamo ogni valore nel proprio percentile:
\[
 U_1=F_1(X_1),\qquad U_2=F_2(X_2).
\]
Se le marginali sono continue, $U_1$ e $U_2$ sono uniformi su $[0,1]$. Un
valore vicino a zero significa ``molto basso rispetto alla propria marginale'';
un valore vicino a uno significa ``molto alto''.

La copula $C$ descrive come questi percentili si combinano. Il teorema di Sklar
scrive la distribuzione congiunta come
\begin{equation}
 F(x_1,x_2)=C\bigl(F_1(x_1),F_2(x_2)\bigr).
 \label{eq:sklar}
\end{equation}

\Idea Le marginali dicono quanto sono rischiosi i due asset separatamente. La
copula dice se i loro giorni buoni e cattivi tendono ad arrivare insieme.
Separare questi due oggetti permette, per esempio, di usare marginali con code
pesanti e una struttura di dipendenza scelta appositamente.

\begin{figure}[ht]
\centering
\begin{tikzpicture}[font=\small,>=Latex]
  \node[flowbox,minimum width=3.3cm] (raw) {rendimenti\\$(x_1,x_2)$};
  \node[flowbox,minimum width=3.3cm,right=1.5cm of raw] (marg) {marginali\\$F_1,F_2$};
  \node[flowbox,minimum width=3.3cm,right=1.5cm of marg] (rank) {percentili\\$(u_1,u_2)$};
  \draw[flowarrow] (raw)--(marg);
  \draw[flowarrow] (marg)--(rank);
  \node[below=0.65cm of marg,align=center] {la copula conserva il legame tra percentili\\e lascia fuori unità di misura e scale marginali};
\end{tikzpicture}

\vspace{0.8em}
\begin{tikzpicture}[x=4.0cm,y=4.0cm,font=\small]
  \draw[axis] (0,0)--(1.08,0) node[right] {$u_1$};
  \draw[axis] (0,0)--(0,1.08) node[above] {$u_2$};
  \draw[gray] (0,0) rectangle (1,1);
  \draw[guide] (0,0)--(1,1) node[above right] {diagonale};
  \fill[mainorange!18] (0,0) rectangle (0.25,0.25);
  \fill[mainblue!12] (0.75,0.75) rectangle (1,1);
  \node[align=center] at (0.13,0.13) {lower\\tail};
  \node[align=center] at (0.87,0.87) {upper\\tail};
  \foreach \x/\y in {0.10/0.14,0.19/0.08,0.31/0.38,0.44/0.40,0.58/0.63,0.68/0.72,0.78/0.76,0.91/0.89}{\fill[data] (\x,\y) circle (1.7pt);}
\end{tikzpicture}
\caption{Prima si trasformano le marginali in percentili; poi si studia la copula nel quadrato unitario. La diagonale confronta la stessa soglia nelle due componenti.}
\label{fig:copula-square}
\end{figure}

\subsection{La sezione diagonale}

La sezione diagonale di una copula bivariata è
\begin{equation}
 \delta_C(u)=C(u,u)=\Prob(U_1\leq u,U_2\leq u).
 \label{eq:diagonal}
\end{equation}
Per esempio, $\delta_C(0.05)$ è la probabilità che entrambe le variabili siano
sotto il proprio quinto percentile. La diagonale non descrive tutta la copula,
ma contiene esattamente l'informazione necessaria per TDC e TCF.

\Attenzione Copule diverse possono avere la stessa diagonale. Durante et al.
studiano anche come costruire copule con diagonale prescritta; questa parte è
importante teoricamente, ma alla prima lettura basta ricordare che la diagonale
è una compressione utile, non l'intera struttura di dipendenza.

\section{Il primo livello: il tail-dependence coefficient}

\subsection{Lower tail}

Sia $X=(X_1,X_2)^\top$ con marginali continue $F_1,F_2$. Il lower TDC è
\begin{equation}
 \lambda_L=\lim_{u\downarrow0}
 \Prob\!\left(X_1\leq F_1^{-1}(u)\mid
 X_2\leq F_2^{-1}(u)\right).
 \label{eq:lambda-lower-cond}
\end{equation}
Nello spazio della copula diventa
\begin{equation}
 \lambda_L=\lim_{u\downarrow0}\frac{C(u,u)}{u}
 =\lim_{u\downarrow0}\frac{\delta_C(u)}{u}.
 \label{eq:lambda-lower-copula}
\end{equation}

\paragraph{Simbolo per simbolo.}
$u$ è una soglia percentile molto bassa; $C(u,u)$ è la probabilità che
entrambi siano sotto la soglia; dividere per $u$ trasforma il conteggio
congiunto in una probabilità condizionata. Il limite spinge la soglia verso
l'angolo $(0,0)$.

\subsection{Upper tail}

Il ragionamento speculare è
\begin{equation}
 \lambda_U=\lim_{u\uparrow1}
 \Prob\!\left(X_1>F_1^{-1}(u)\mid X_2>F_2^{-1}(u)\right),
 \label{eq:lambda-upper-cond}
\end{equation}
e, tramite la diagonale,
\begin{equation}
 \lambda_U=\lim_{u\uparrow1}
 \frac{1-2u+C(u,u)}{1-u}.
 \label{eq:lambda-upper-copula}
\end{equation}
Il numeratore $1-2u+C(u,u)$ è la probabilità che entrambi superino $u$.

\begin{figure}[ht]
\centering
\begin{tikzpicture}[x=4.0cm,y=4.0cm,font=\small]
  \draw[axis] (0,0)--(1.08,0) node[right] {$u_1$};
  \draw[axis] (0,0)--(0,1.08) node[above] {$u_2$};
  \draw[gray] (0,0) rectangle (1,1);
  \draw[guide] (0.18,0)--(0.18,1);
  \draw[guide] (0,0.18)--(1,0.18);
  \fill[mainred!18] (0,0) rectangle (0.18,0.18);
  \node[align=center] at (0.09,0.09) {$C(u,u)$};
  \draw[-{Latex},mainred,thick] (0.42,0.35)--(0.20,0.20);
  \node[align=left] at (0.67,0.38) {spingere $u$ verso $0$\\significa entrare nell'angolo};
  \node[align=center] at (0.52,-0.18) {$\lambda_L$ osserva il rapporto tra massa congiunta e massa marginale};
\end{tikzpicture}
\caption{Il TDC è una misura di bordo: osserva che cosa resta quando la soglia raggiunge l'angolo estremo.}
\label{fig:tdc-corner}
\end{figure}

\subsection{Dipendenza e indipendenza asintotica}

Se $\lambda_L>0$, la lower tail è asintoticamente dipendente; se
$\lambda_L=0$, è asintoticamente indipendente. Lo stesso vale per la upper
tail.

La normale bivariata ha TDC nullo in entrambe le code per ogni
$\rho\in(-1,1)$. La Student-$t$ bivariata ha invece TDC positivo: le code
pesanti mantengono una probabilità limite di estremi congiunti.

\Ricorda $\lambda=0$ non significa indipendenza a ogni soglia. Significa che
la probabilità condizionata tende a zero nel limite. A una soglia finita può
essere ancora economicamente rilevante. Questo è il punto da cui parte
Hua--Joe.

\section{Il secondo livello: il tail order e la velocità di decadimento}

\subsection{Due zeri possono essere diversi}

Supponiamo che due probabilità congiunte abbiano, in prima approssimazione, le
forme $u^{3/2}$ e $u^2$. Entrambe tendono a zero quando $u\downarrow0$, quindi
il TDC può essere nullo in entrambi i casi. Però, con $u=0.01$:
\[
 u^{3/2}=10^{-3},\qquad u^2=10^{-4}.
\]
La prima probabilità è dieci volte la seconda.

\begin{figure}[ht]
\centering
\begin{tikzpicture}[x=1.15cm,y=0.85cm,font=\small]
  \draw[axis] (0,0)--(7.1,0) node[right] {$u$ si avvicina a $0$};
  \draw[axis] (0,0)--(0,4.2) node[above] {probabilità congiunta};
  \draw[mainblue,very thick] (0.35,3.55)..controls(1.5,2.75)and(3.7,1.35)..(6.7,0.35);
  \draw[mainred,very thick] (0.35,3.55)..controls(1.3,1.75)and(3.4,0.65)..(6.7,0.06);
  \node[mainblue,align=left] at (4.6,2.25) {$u^{3/2}$: scende\\più lentamente};
  \node[mainred,align=left] at (4.3,0.65) {$u^2$: scende\\più rapidamente};
  \node[align=center] at (3.5,-0.68) {entrambe arrivano a zero, ma non alla stessa velocità};
\end{tikzpicture}
\caption{Il tail order distingue velocità che il solo limite nullo non separa. Diagramma concettuale.}
\label{fig:orders-curves}
\end{figure}

\subsection{La formula del tail order}

Hua--Joe scrivono la diagonale lower come
\begin{equation}
 C(u\one)\sim u^{\kappa_L(C)}\ell_L(u),
 \qquad u\downarrow0,
 \label{eq:tail-order-lower}
\end{equation}
e la survival copula upper come
\begin{equation}
 \overline C((1-u)\one)\sim u^{\kappa_U(C)}\ell_U(u).
 \label{eq:tail-order-upper}
\end{equation}

Qui $\kappa$ è l'esponente che controlla la velocità principale;
$\ell(u)$ è una funzione lentamente variabile, per esempio una potenza di
$-\log u$; il simbolo $\sim$ significa che il rapporto tra i due membri tende
a uno.

\begin{figure}[ht]
\centering
\begin{tikzpicture}[font=\small]
  \draw[very thick,mainblue] (0,0)--(12,0);
  \foreach \x in {0,6,12}{\draw[thick] (\x,-0.15)--(\x,0.15);}
  \node[below=0.22cm] at (0,0) {$\kappa=1$};
  \node[below=0.22cm] at (6,0) {$1<\kappa<d$};
  \node[below=0.22cm] at (12,0) {$\kappa=d$};
  \node[above=0.28cm,align=center] at (0,0) {tail dependence\\classica};
  \node[above=0.28cm,align=center] at (6,0) {intermediate\\tail dependence};
  \node[above=0.28cm,align=center] at (12,0) {indipendenza\\come riferimento};
  \draw[-{Latex},mainred,thick] (4.5,1.7)--(1.0,1.7);
  \node[above] at (2.8,1.7) {ordine più piccolo $=$ decadimento più lento};
\end{tikzpicture}
\caption{La scala del tail order in dimensione $d$, sotto condizioni di dipendenza positiva.}
\label{fig:kappa-scale}
\end{figure}

\paragraph{I tre casi.}
Se $\kappa=1$ e $\ell(u)$ tende a una costante positiva, ritroviamo la tail
dependence classica. Se $1<\kappa<d$, il TDC è nullo ma la probabilità scende
più lentamente dell'indipendenza: è intermediate tail dependence. Se
$\kappa=d$, abbiamo la velocità dell'indipendenza, perché
$C(u,\ldots,u)=u^d$ per la copula indipendente.

\Attenzione Hua--Joe usano anche il simbolo $\lambda(C)$ per il limite di
$\ell(u)$. Non va confuso con il TDC di Schmidt. I due coincidono nel caso
$\kappa=1$ con limite positivo; quando $\kappa>1$, il TDC classico è zero ma
$\kappa$ resta informativo.

\subsection{L'esempio gaussiano}

Per una copula gaussiana bivariata con correlazione $\rho>0$:
\begin{equation}
 \kappa=\frac{2}{1+\rho},
 \qquad
 \ell(u)=(-\log u)^{-\rho/(1+\rho)}.
 \label{eq:gaussian-order}
\end{equation}
Con $\rho=0.5$, $\kappa=4/3$: il TDC è ancora zero, ma l'ordine è minore di
$2$. La correlazione positiva rallenta quindi il decadimento rispetto
all'indipendenza completa.

\subsection{Percorsi non diagonali}

Finora abbiamo fatto entrare tutte le variabili nella coda usando la stessa
soglia: ci siamo mossi lungo la diagonale. Hua--Joe osservano che non è
l'unico percorso possibile. Un titolo potrebbe essere nel proprio 1\% peggiore
mentre un altro è nel 5\% peggiore. La \emph{tail order function} serve a
descrivere anche questi avvicinamenti sbilanciati all'angolo.

Per questa prima visione d'insieme non serve la formula completa. Basta
ricordare che due copule con lo stesso $\kappa$ lungo la diagonale possono
ancora differire quando le soglie delle variabili non sono uguali.

\Salta alla prima lettura le prove con variazione regolare, trasformate di
Laplace e derivate di ordine elevato. Devi ricordare la direzione: $\kappa$ più
piccolo significa coda più persistente.

\section{Il terzo livello: la tail concentration function a soglia finita}

\subsection{Dal limite alla curva}

Durante et al. prendono le stesse frazioni dei TDC ma non mandano subito la
soglia al limite. Definiscono
\begin{equation}
 q_C(t)=
 \begin{cases}
  \displaystyle\frac{C(t,t)}{t}, & 0<t\leq0.5,\\[1.1ex]
  \displaystyle\frac{1-2t+C(t,t)}{1-t}, & 0.5<t<1.
 \end{cases}
 \label{eq:tcf}
\end{equation}

Se $t<0.5$, la curva misura la concentrazione nella lower tail; se $t>0.5$,
misura la upper tail. Dividiamo per $t$ o $1-t$ per confrontare la massa
congiunta con la probabilità marginale della stessa soglia.

\Esempio Se $t=0.05$ e la probabilità congiunta osservata è $0.02$, allora
$q_C(0.05)=0.02/0.05=0.4$. Tra le occasioni associate al quinto percentile
marginale, la quota congiunta normalizzata è $40\%$.

I limiti della curva recuperano i TDC:
\[
 \lim_{t\downarrow0}q_C(t)=\lambda_L,
 \qquad
 \lim_{t\uparrow1}q_C(t)=\lambda_U.
\]

\begin{figure}[ht]
\centering
\begin{tikzpicture}[x=1.02cm,y=0.78cm,font=\small]
  \draw[axis] (0,0)--(10.3,0) node[right] {$t$};
  \draw[axis] (0,0)--(0,4.3) node[above] {$q_C(t)$};
  \draw[guide] (5,0)--(5,4) node[above] {$0.5$};
  \draw[mainblue,very thick] (0.3,0.45)..controls(1.7,0.6)and(3.5,1.25)..(5,2.0)
    ..controls(6.4,2.35)and(8.6,2.75)..(9.8,3.0);
  \draw[mainred,very thick] (0.3,0.12)..controls(1.7,0.25)and(3.5,0.75)..(5,2.0)
    ..controls(6.4,2.95)and(8.6,3.55)..(9.8,3.88);
  \node[mainblue,align=center] at (2.2,2.15) {curva più alta\\$=$ più concentrazione};
  \node[mainred,align=center] at (7.4,1.35) {convergenza lenta\\$=$ rischio finito persistente};
  \node[align=center] at (5,-0.75) {lower tail $\longleftarrow$ centro $\longrightarrow$ upper tail};
\end{tikzpicture}
\caption{La TCF è una curva: altezza e velocità di convergenza contengono informazione che il solo valore limite perde.}
\label{fig:tcf-reading}
\end{figure}

Nel punto centrale
\begin{equation}
 q_C(0.5)=\frac{1+\beta_C}{2},
 \label{eq:blomqvist}
\end{equation}
dove $\beta_C$ è la misura di associazione di Blomqvist. Il paper confronta
copule con lo stesso $\beta_C$: stesso livello centrale, curve di coda diverse.

\subsection{Come leggere alcune famiglie}

La copula Gaussiana ha TDC nullo, ma può convergere lentamente al bordo.
Frank e Plackett hanno anch'esse TDC nullo e, nel confronto del paper, code che
si spengono più rapidamente.

Clayton concentra dipendenza nella lower tail: è spesso una scelta conservativa
se valori bassi rappresentano perdite. Survival Clayton riflette il fenomeno
nella upper tail. Gumbel e Galambos privilegiano la upper tail e mostrano TCF
molto vicine nel confronto. La Student-$t$ presenta dipendenza positiva in
entrambe le code, coerente con code pesanti e simmetria.

\Ricorda La famiglia non cambia i dati: è un'ipotesi sulla forma del legame.
Famiglie con associazione centrale simile possono implicare rischi estremi
molto diversi.

\section{Una mappa unica: TDC, tail order e TCF}

Le tre misure non sono concorrenti. Rispondono a tre domande successive.

\begin{figure}[ht]
\centering
\begin{tikzpicture}[font=\small,node distance=0.75cm]
  \node[flowbox,text width=3.35cm] (tdc) {TDC $\lambda$\\il limite è positivo o nullo?};
  \node[flowbox,text width=3.35cm,right=of tdc] (kap) {tail order $\kappa$\\se è nullo, quanto velocemente?};
  \node[flowbox,text width=3.35cm,right=of kap] (tcf) {TCF $q_C(t)$\\che cosa accade a soglie finite?};
  \draw[flowarrow] (tdc)--(kap);
  \draw[flowarrow] (kap)--(tcf);
  \node[below=0.9cm of kap,align=center] {stesso oggetto di fondo: probabilità congiunta vicino agli angoli della copula};
\end{tikzpicture}
\caption{Mappa concettuale delle tre lenti: limite, velocità, percorso finito.}
\label{fig:three-measures}
\end{figure}

\begin{center}
\begin{tabularx}{\textwidth}{>{\bfseries}p{2.4cm}X X}
\toprule
Oggetto & Domanda & Limite principale\\
\midrule
TDC $\lambda$ & rimane una probabilità positiva all'estremo? & comprime tutto in un numero e raggruppa molti casi sotto $0$\\
Tail order $\kappa$ & con quale potenza scende la probabilità? & è asintotico e richiede condizioni di regolarità\\
TCF $q_C(t)$ & quanto è concentrata la coda alla soglia $t$? & dipende dalla soglia e, empiricamente, dalla numerosità campionaria\\
\bottomrule
\end{tabularx}
\end{center}

\subsection*{Un unico esempio letto con tutte e tre le lenti}

Confrontiamo due modelli ipotetici nella lower tail:
\[
 C_A(u,u)\sim u^{3/2},
 \qquad
 C_B(u,u)\sim u^2.
\]
Non stiamo dicendo che queste formule descrivano tutto il modello. Le usiamo
per osservare, una dopo l'altra, le tre domande centrali.

\paragraph{Prima domanda: che cosa dice il TDC?}
Dividiamo la probabilità congiunta per la probabilità marginale $u$:
\[
 \frac{C_A(u,u)}{u}\sim u^{1/2},
 \qquad
 \frac{C_B(u,u)}{u}\sim u.
\]
Entrambi i rapporti tendono a zero. Il TDC lower è quindi $\lambda_L=0$ per
entrambi: il primo strumento li colloca nella stessa categoria di indipendenza
asintotica.

\paragraph{Seconda domanda: che cosa aggiunge il tail order?}
Nel modello $A$ l'esponente è $\kappa_A=3/2$; nel modello $B$ è
$\kappa_B=2$. Il modello $A$ scende più lentamente e conserva più probabilità
di estremi congiunti prima del limite. Se $B$ rappresenta la velocità
dell'indipendenza bivariata, $A$ si trova nella regione intermedia.

\paragraph{Terza domanda: che cosa vediamo a una soglia finita?}
Prendiamo $u=0.05$. Otteniamo:
\[
 C_A(0.05,0.05)\approx0.05^{3/2}\approx0.0112,
 \qquad
 C_B(0.05,0.05)=0.05^2=0.0025.
\]
Normalizzando come nella parte lower della TCF:
\[
 q_A(0.05)\approx\frac{0.0112}{0.05}=0.224,
 \qquad
 q_B(0.05)=\frac{0.0025}{0.05}=0.05.
\]
A questa soglia il modello $A$ produce una concentrazione congiunta più di
quattro volte quella del modello $B$, pur avendo lo stesso TDC nullo.

\Ricorda Il TDC assegna la categoria limite; $\kappa$ distingue la velocità;
la TCF mostra quanto la differenza pesa alla soglia che possiamo davvero
osservare.

\section{Dalla teoria alla stima}

\subsection{Stimare il TDC con ranghi estremi}

Siano $m$ osservazioni bivariate e $R_1^{(j)},R_2^{(j)}$ i loro ranghi. Si
sceglie un numero $k$ di osservazioni estreme. Il primo stimatore upper di
Schmidt conta le osservazioni che appartengono contemporaneamente alle due
code:
\begin{equation}
 \widehat\lambda_{U,m}^{(1)}=\frac{1}{k}\sum_{j=1}^m
 \ind\!\left(R_1^{(j)}>m-k,\ R_2^{(j)}>m-k\right).
 \label{eq:tdc-estimator}
\end{equation}

Con $m=20$ e $k=4$, i ranghi estremi sono $17,18,19,20$. Se due osservazioni
sono estreme in entrambe le componenti, la stima è $2/4=0.5$.

Il valore di $k$ crea un compromesso: $k$ piccolo significa pochi punti e
grande varianza; $k$ grande include osservazioni meno estreme e aumenta il
bias. In teoria $k\to\infty$ ma $k/m\to0$.

Schmidt considera anche un secondo stimatore non parametrico basato
sull'unione delle due code e uno parametrico per distribuzioni ellittiche,
$\widehat\lambda^{(3)}=\lambda(\widehat\alpha,\widehat\rho)$. Nelle sue
simulazioni $t$ bivariate, il parametrico funziona meglio quando il modello
ellittico è corretto; il non parametrico è più prudente quando la forma della
distribuzione è ignota.

\subsection{Stimare la TCF}

Con pseudo-osservazioni $(U_i,V_i)$, la copula empirica è
\begin{equation}
 C_n(u,v)=\frac1n\sum_{i=1}^n
 \ind\!\left(U_i\leq u,V_i\leq v\right).
 \label{eq:empirical-copula}
\end{equation}
Sostituendo $C_n(t,t)$ nella formula della TCF otteniamo
\begin{equation}
 q_{\mathrm{emp}}(t)=
 \frac{C_n(t,t)}{t}\ind_{(0,0.5]}(t)+
 \frac{1-2t+C_n(t,t)}{1-t}\ind_{(0.5,1)}(t).
 \label{eq:empirical-tcf}
\end{equation}

Vicino a zero il quadrato $[0,t]^2$ contiene pochi punti. La curva può quindi
saltare o diventare rumorosa. Durante et al. consigliano di iniziare da una
soglia $\varepsilon>0$ legata alla dimensione del campione, invece di leggere
la curva fino al bordo come se fosse precisa.

\Attenzione Una TCF empirica è uno strumento diagnostico. Per fare inferenza
servono anche incertezza, sensibilità alla soglia e confronto con metodi di
extreme-value theory.

\section{Famiglie di modelli: che cosa stiamo scegliendo?}

\subsection{Copule archimedee}

Le copule archimedee sono famiglie costruite a partire da una sola funzione,
chiamata \emph{generatore}. Il generatore funziona come una regola che decide
quanto e in quale parte del quadrato i percentili tendono ad avvicinarsi.

Per ora non serve ricordare la formula del generatore. È più utile distinguere
i comportamenti: Clayton concentra la dipendenza nella lower tail;
Gumbel--Hougaard nella upper tail; Frank e Ali--Mikhail--Haq possono descrivere
dipendenza centrale senza un TDC positivo.

Hua--Joe mostrano che, in alcune costruzioni, il generatore può essere
interpretato attraverso un fattore comune non osservato. L'intuizione è che
uno stesso shock latente può spingere più variabili verso un bordo. Le
trasformate di Laplace usate per rendere rigorosa questa idea possono essere
saltate alla prima lettura.

\subsection{Distribuzioni ellittiche}

Le distribuzioni ellittiche producono nuvole di punti con contorni simili a
ellissi. La normale e la Student-$t$ appartengono entrambe a questa classe, ma
non hanno lo stesso comportamento estremo: la normale ha code sottili e TDC
nullo, mentre la Student-$t$ ha code pesanti e TDC positivo.

La forma ellittica è relativamente rigida e simmetrica. Se i dati mostrano
perdite congiunte molto più forti dei guadagni congiunti, una copula
asimmetrica può essere più appropriata.

\subsection{Come confrontare le famiglie}

Non confrontiamo direttamente parametri con nomi simili tra famiglie. Prima
calibriamo un livello di associazione comparabile, poi osserviamo:
\begin{itemize}
  \item quale coda può essere dipendente;
  \item se la dipendenza è simmetrica;
  \item quali $\lambda$, $\kappa$ e TCF produce;
  \item se l'ipotesi di forma è plausibile per i dati.
\end{itemize}

\section{Applicazioni finanziarie: dal VaR al clustering}

\subsection{VaR e dipendenza di coda}

Il VaR è un quantile elevato della perdita. Per un portafoglio dipende da due
componenti: le code marginali dei singoli asset e la probabilità che le perdite
arrivino insieme.

Schmidt confronta una normale bivariata, una $t$ bivariata e una copula-$t$
con marginali-$t$. La normale tende a sottostimare quantili estremi per via
delle code sottili. La $t$ bivariata funziona bene quando la struttura è
ellittica. Se la copula e le marginali richiedono gradi di libertà diversi,
separarle con una copula-$t$ offre più flessibilità.

Nei dati BMW--Deutsche Bank Schmidt stima un upper TDC circa $0.31$; nei dati
DEM/USD--JPY/USD circa $0.17$, scegliendo $k$ in una regione di plateau. Questi
numeri sono esempi empirici, non valori universali.

\subsection{Scegliere una copula con una mappa di distanze}

Durante et al. confrontano la TCF empirica con le TCF di copule candidate:
\begin{equation}
 d(C_n,C_i)=\int_0^1
 \bigl(q_{\mathrm{emp}}(t)-q_{C_i}(t)\bigr)^2\,dt.
 \label{eq:tcf-distance}
\end{equation}
Le distanze tra tutte le curve vengono rappresentate in un piano tramite
Multidimensional Scaling (MDS). È come costruire una mappa conoscendo le
distanze tra città: punti vicini hanno TCF simili. Il piano è una
visualizzazione approssimata, non una nuova variabile economica.

\subsection{Quali mercati perdono insieme?}

Per raggruppare serie finanziarie, il paper procede dalla dinamica marginale
alla dipendenza di coda. Ogni serie viene filtrata con un modello ARMA--GARCH
con innovazioni Student-$t$; i residui vengono trasformati in percentili tramite
PIT; per ogni coppia si calcola la TCF lower; le distanze entrano in un
clustering gerarchico.

\begin{figure}[ht]
\centering
\begin{tikzpicture}[font=\small,node distance=0.5cm]
  \node[flowbox,minimum width=2.35cm] (ret) {rendimenti\\grezzi};
  \node[flowbox,minimum width=2.55cm,right=of ret] (filter) {ARMA--GARCH\\marginali};
  \node[flowbox,minimum width=2.35cm,right=of filter] (pit) {residui\\$\to$ PIT};
  \node[flowbox,minimum width=2.35cm,right=of pit] (tcf) {TCF lower\\a coppie};
  \node[flowbox,minimum width=2.35cm,right=of tcf] (cluster) {matrice\\$\to$ cluster};
  \draw[flowarrow] (ret)--(filter);
  \draw[flowarrow] (filter)--(pit);
  \draw[flowarrow] (pit)--(tcf);
  \draw[flowarrow] (tcf)--(cluster);
  \node[below=0.8cm of filter,align=center] {qui si tolgono\\scala e volatilità};
  \node[below=0.8cm of tcf,align=center] {qui si misura chi\\perde insieme};
\end{tikzpicture}
\caption{Pipeline finanziaria: i modelli marginali preparano serie confrontabili; la copula e la TCF misurano la dipendenza tra le perdite.}
\label{fig:financial-pipeline}
\end{figure}

Nel dataset di 23 indici MSCI sviluppati (2002--2010), il complete-linkage
clustering suggerisce quattro gruppi, interpretabili soprattutto in senso
geografico. Una possibile idea di diversificazione è scegliere asset da
cluster diversi, evitando di concentrare il portafoglio in mercati che perdono
insieme. È un criterio esplorativo, non una strategia completa.

\section{Come affrontare un'analisi empirica senza perdersi}

Una sequenza prudente potrebbe essere:
\begin{enumerate}
  \item chiarire se valori bassi o alti rappresentano il rischio rilevante;
  \item modellare o trasformare bene le marginali;
  \item osservare scatterplot di percentili e copula empirica;
  \item tracciare la TCF con una regione di soglie supportata dal campione;
  \item stimare TDC e, se utile, tail order con analisi di sensibilità;
  \item confrontare famiglie con code differenti, non solo fit centrale;
  \item valutare come la scelta cambia VaR, clustering o altra misura finale;
  \item riportare incertezza, soglie e ipotesi.
\end{enumerate}

\Ricorda Una bella curva non elimina il problema statistico. I risultati
dipendono da numerosità, non-stazionarietà, scelta della soglia, modello
marginale e famiglia della copula.

\newpage
\section{Sintesi finale}

\subsection{La storia completa in dieci passaggi}

\begin{enumerate}
  \item I rischi estremi congiunti non sono descritti completamente dalla
  correlazione.
  \item Le marginali descrivono ogni variabile; la copula descrive il legame tra
  percentili.
  \item La diagonale $C(u,u)$ guarda soglie uguali e conduce direttamente ai
  coefficienti di coda.
  \item Il TDC $\lambda$ chiede se una probabilità condizionata estrema resta
  positiva nel limite.
  \item La normale ha TDC nullo; la Student-$t$ ha TDC positivo.
  \item Un TDC nullo contiene casi molto diversi: il tail order $\kappa$ misura
  la velocità di decadimento.
  \item $\kappa=1$ corrisponde alla dipendenza classica, $1<\kappa<d$ alla
  regione intermedia, $\kappa=d$ alla velocità di indipendenza.
  \item La TCF $q_C(t)$ mostra lower e upper tail a soglie finite e recupera i
  TDC ai bordi.
  \item Stima di TDC e TCF significa contare ranghi estremi: pochi punti
  producono grande incertezza.
  \item In finanza queste idee influenzano VaR, scelta della copula e
  diversificazione tra mercati che perdono insieme.
\end{enumerate}

\subsection{Confronto dei tre paper}

\begin{center}
\small
\begin{tabularx}{\textwidth}{>{\bfseries}p{2.4cm}p{3.1cm}X X}
\toprule
Paper & Domanda centrale & Contributo da ricordare & Utilità per noi\\
\midrule
Schmidt (2005) & Come definire e stimare la dipendenza estrema? & TDC, copule, stimatori, simulazione e VaR & base introduttiva e collegamento diretto con finanza\\
Hua--Joe (2011) & Come distinguere i casi con TDC nullo? & tail order $\kappa$, funzione di ordine, intermediate dependence, asimmetria & raffina la categoria di indipendenza asintotica\\
Durante et al. (2015) & Come leggere e usare la dipendenza a soglie finite? & diagonale, TCF, confronto di copule, MDS e clustering & porta il discorso verso diagnostica e applicazioni empiriche\\
\bottomrule
\end{tabularx}
\end{center}

\subsection{Mappa concettuale}

\begin{center}
\begin{tikzpicture}[font=\small,node distance=1.0cm and 1.2cm]
  \node[flowbox] (data) {dati\\$(X_1,X_2)$};
  \node[flowbox,right=of data] (marg) {marginali\\$F_1,F_2$};
  \node[flowbox,right=of marg] (cop) {copula\\$C$};
  \node[flowbox,below=of cop] (diag) {diagonale\\$C(u,u)$};
  \node[flowbox,left=of diag] (tdc) {TDC\\$\lambda$};
  \node[flowbox,left=of tdc] (kap) {tail order\\$\kappa$};
  \node[flowbox,right=of diag] (tcf) {TCF\\$q_C(t)$};
  \node[flowbox,below=of diag] (risk) {VaR, scelta copula,\\clustering, diversificazione};
  \draw[flowarrow] (data)--(marg);
  \draw[flowarrow] (marg)--(cop);
  \draw[flowarrow] (cop)--(diag);
  \draw[flowarrow] (diag)--(tdc);
  \draw[flowarrow] (tdc)--(kap);
  \draw[flowarrow] (diag)--(tcf);
  \draw[flowarrow] (tdc)--(risk);
  \draw[flowarrow] (kap)--(risk);
  \draw[flowarrow] (tcf)--(risk);
\end{tikzpicture}
\end{center}

\subsection{Glossario e pronuncia}

\begin{description}[leftmargin=4.0cm,style=nextline]
  \item[$X_i$] variabile originale, per esempio rendimento o perdita dell'asset
  $i$.
  \item[$F_i$] distribuzione marginale della variabile $i$ presa da sola.
  \item[$U_i=F_i(X_i)$] percentile della variabile $i$, uniforme tra zero e uno.
  \item[$C$] copula, cioè distribuzione congiunta dei percentili; si legge
  ``ci''.
  \item[$\delta_C(u)$] diagonal section $C(u,u)$; delta si legge ``delta''.
  \item[$\overline C$] survival copula, usata per probabilità di superamento
  congiunto.
  \item[$\lambda_L,\lambda_U$] lower e upper TDC di Schmidt; lambda si legge
  ``lambda''.
  \item[$\kappa_L,\kappa_U$] tail order di Hua--Joe; kappa è la lettera greca
  $\kappa$ e si legge ``cappa''.
  \item[$k$] numero o indice di osservazioni estreme negli stimatori; è lettera
  latina ``k'', letta spesso ``cappa'', ma non è $\kappa$.
  \item[$\ell(u)$] funzione lentamente variabile che corregge la potenza
  $u^\kappa$.
  \item[$b(w)$] tail order function lungo percorsi non diagonali.
  \item[$q_C(t)$] TCF, concentrazione di coda alla soglia finita $t$; si legge
  ``cu di ti''.
  \item[$\beta_C$] associazione centrale di Blomqvist; beta si legge ``beta''.
  \item[$\rho$] correlazione o parametro di dipendenza ellittico; rho si legge
  ``ro''.
  \item[$\alpha$] indice di coda o parametro di famiglia, a seconda del
  contesto; alpha si legge ``alfa''.
  \item[$\theta$] parametro generico di una famiglia di copule; theta si legge
  ``teta''.
  \item[$\nu$] gradi di libertà della distribuzione Student-$t$; nu si legge
  ``ni'' oppure ``nu''.
  \item[$C_n$] copula empirica ottenuta dai ranghi osservati.
  \item[PIT] probability integral transform, cioè trasformazione di un residuo
  nel proprio percentile stimato.
  \item[MDS] Multidimensional Scaling, mappa approssimata costruita da una
  matrice di distanze.
  \item[ARMA--GARCH] modello marginale per dinamica media e volatilità; non è
  la copula.
\end{description}

\subsection{Domande concrete per il relatore}

\paragraph{Concettuali.}
\begin{itemize}
  \item Nella tesi il focus deve essere TDC, tail order, TCF o il confronto tra
  i tre?
  \item Per rendimenti finanziari è meglio parlare di lower tail dei rendimenti
  o upper tail delle perdite?
  \item Quanto è importante distinguere dipendenza asintotica e dipendenza a
  soglia finita?
\end{itemize}

\paragraph{Metodologiche.}
\begin{itemize}
  \item Come scegliere $k$ per il TDC e $\varepsilon$ per la TCF empirica?
  \item Quale stimatore di $\kappa$ è più adatto a campioni finanziari?
  \item Come costruire intervalli o bande di confidenza per TDC e TCF?
  \item Conviene imporre una copula parametrica oppure partire da ranghi e
  strumenti non parametrici?
\end{itemize}

\paragraph{Empiriche.}
\begin{itemize}
  \item Quali dati dovremmo usare: indici, banche, valute, crypto o portafogli?
  \item Quale frequenza e quale orizzonte temporale sono realistici?
  \item Dobbiamo filtrare i rendimenti con ARMA--GARCH prima di studiare la
  copula?
  \item Il risultato finale dovrebbe essere VaR, confronto di copule,
  clustering o simulazione metodologica?
\end{itemize}

\subsection{Possibili direzioni di tesi, senza scegliere ancora}

\begin{enumerate}
  \item \emph{Confronto metodologico}: TDC, tail order e TCF sulle stesse
  famiglie e sugli stessi dati.
  \item \emph{Asimmetria finanziaria}: confronto tra lower e upper tail per
  indici o settori.
  \item \emph{Rischio di portafoglio}: sensibilità del VaR alla copula e alla
  dipendenza estrema.
  \item \emph{Clustering}: raggruppare asset per perdite congiunte e confrontare
  il risultato con cluster basati su correlazione.
  \item \emph{Simulazione}: studiare bias, varianza e scelta della soglia per
  stimatori di dipendenza di coda.
\end{enumerate}

Queste sono direzioni possibili, non una proposta definitiva. La scelta dipende
da ciò che il relatore considera più interessante e dal livello matematico ed
empirico atteso.

\subsection{Parti da rileggere}

\begin{itemize}
  \item Schmidt, Sec.~3.2--3.4: definizioni, copule e stimatori del TDC.
  \item Schmidt, Sec.~3.5--3.7: simulazione, rendimenti e VaR.
  \item Hua--Joe, Definizioni 1--3 ed esempio gaussiano: $\kappa$, $\ell$ e
  intermediate dependence.
  \par\item Hua--Joe, Sec.~3: interpretazione archimedea tramite trasformata di
  Laplace, senza concentrarsi subito sulle prove.
  \par\item Durante et al., Sec.~3: TCF teorica ed empirica.
  \par\item Durante et al., Sec.~4: mappa MDS e clustering finanziario.
\end{itemize}

\subsection{Riferimenti essenziali}

\begin{itemize}
  \item R.~Schmidt (2005), \emph{Tail Dependence}, capitolo introduttivo su
  definizione, calcolo, stima e VaR.
  \item L.~Hua e H.~Joe (2011), \emph{Tail order and intermediate tail
  dependence of multivariate copulas}, Journal of Multivariate Analysis.
  \item F.~Durante, J.~Fernández-Sánchez e R.~Pappadà (2015), \emph{Copulas,
  diagonals, and tail dependence}, Fuzzy Sets and Systems.
  \item R.~B.~Nelsen, \emph{An Introduction to Copulas}, per le basi sulle
  copule.
  \item A.~J.~McNeil, R.~Frey e P.~Embrechts, \emph{Quantitative Risk
  Management}, per il collegamento con finanza e rischio.
\end{itemize}

\paragraph{Ultima frase da tenere a mente.}
La tail dependence non è un unico numero: è un insieme di domande sul limite,
sulla velocità e sulla scala finita. TDC, tail order e TCF sono tre modi
complementari di rendere queste domande precise.

\end{document}
